% ==============================================================================
% CURRICULUM VITAE LATEX TEMPLATE & AGENT GENERATION INSTRUCTIONS
% ==============================================================================
% INSTRUCTIONS FOR THE AI AGENT:
% 1. PURPOSE:
%    This document serves as the master LaTeX template for generating customized,
%    ATS-optimized CVs tailored to specific job descriptions.
%
% 2. SOURCE OF TRUTH:
%    Populate this document ONLY with verified career data from:
%    `EN/CareerMasterDocument/`
%    - Personal_Information.md : Contact details, links, location.
%    - Professional_Summary.md : Base summaries and value propositions.
%    - Work_Experience/        : Chronological roles, responsibilities, impact.
%    - Projects.md             : Featured technical and business projects.
%    - Education.md            : Degrees, universities, dates, coursework.
%    - Certifications.md       : Professional credentials and licensures.
%    - Skills.md & Technologies.md: Technical competencies and tools.
%    - Keywords.md             : Industry and ATS keywords.
%    - Quantifiable_Results.md : Metrics, KPIs, and concrete outcomes.
%
% 3. IMMUTABLE FACTS RULE:
%    - NEVER alter, estimate, round, or fabricate employment dates.
%    - NEVER change company names, job titles, or educational institutions.
%    - Maintain exact reverse chronological order.
%
% 4. TAILORING & SELECTION:
%    - Analyze the target job description for key requirements, tools, and keywords.
%    - Select the most relevant bullets and projects. Do not include everything.
%    - Emphasize quantifiable business outcomes (STAR method) and bold key technologies.
%
% 5. LATEX SYNTAX & ESCAPING:
%    - ALWAYS escape reserved LaTeX characters:
%      '&' -> '\&', '%' -> '\%', '_' -> '\_', '#' -> '\#', '$' -> '\$'
%    - Preserve all custom macros (\sectionline, \cvblock, \jobentry) and layout dimensions.
%    - Maintain document length (typically strictly 2 pages) by balancing bullet counts.
% ==============================================================================

\RequirePackage{cmap} % Prevents accent separation when copying PDF text (ATS Optimization)
\documentclass[10pt,a4paper]{article}

% Packages
\usepackage[T1]{fontenc}
\usepackage[utf8]{inputenc}
\usepackage[margin=1.2cm, top=0.8cm, bottom=1.2cm]{geometry}
\usepackage{titlesec}
\usepackage{enumitem}
\usepackage{hyperref}
\usepackage{setspace}
\usepackage{xcolor}
\usepackage{tcolorbox}
\usepackage{helvet}
\renewcommand{\familydefault}{\sfdefault}

% Colors matching reference design
\definecolor{bannerblue}{HTML}{BBD0E3}
\definecolor{dividerblue}{HTML}{8BA5BE}
\definecolor{darkblue}{HTML}{2C4A6F}
\definecolor{textdark}{HTML}{222222}

% PDF Metadata
% INSTRUCTION: Replace bracketed placeholders with candidate name and target role
\hypersetup{
    colorlinks=true,
    urlcolor=darkblue,
    linkcolor=darkblue,
    pdfauthor={[INSTRUCTION: Insert Candidate Full Name]},
    pdftitle={[INSTRUCTION: Insert 'CV - Candidate Full Name - Target Role Title']}
}

% Global Spacing & List Formatting
\setstretch{1.04}
\setlist[itemize]{leftmargin=12pt, itemsep=2pt, topsep=1pt, parsep=0pt, partopsep=0pt}

% Section divider line macro (thick block on left, thin line on right)
\newcommand{\sectionline}{%
  \vspace{4pt}%
  \noindent%
  \begin{minipage}[b]{3.5cm}%
    \color{bannerblue}\rule{\linewidth}{4.5pt}%
  \end{minipage}%
  \hspace{0.3cm}%
  \begin{minipage}[b]{\dimexpr\textwidth-3.8cm\relax}%
    \color{dividerblue}\rule{\linewidth}{1.2pt}%
  \end{minipage}%
  \par \vspace{4pt}%
}

% Robust 2-column block macro with strict width math
% Parameter #1: Section title (left column, 3.5cm)
% Parameter #2: Section content (right column, remaining width)
\newcommand{\cvblock}[2]{%
  \vspace{2pt}%
  \noindent%
  \begin{minipage}[t]{3.5cm}%
    \raggedright \bfseries \small \color{textdark} #1%
  \end{minipage}%
  \hspace{0.3cm}%
  \begin{minipage}[t]{\dimexpr\textwidth-3.8cm\relax}%
    \small \color{textdark} #2%
  \end{minipage}%
  \par \vspace{3pt}%
}

% Macro for subsequent entries sharing the section alignment without repeating the left header
\newcommand{\jobentry}[1]{%
  \vspace{2pt}%
  \noindent%
  \begin{minipage}[t]{3.5cm}%
    \hspace{1pt}%
  \end{minipage}%
  \hspace{0.3cm}%
  \begin{minipage}[t]{\dimexpr\textwidth-3.8cm\relax}%
    \small \color{textdark} #1%
  \end{minipage}%
  \par \vspace{3pt}%
}

\begin{document}
\color{textdark}

%==================================================
% TOP CONTACT BANNER
%==================================================
% INSTRUCTION FOR AGENT:
% Extract candidate location, phone number, and primary email from Personal_Information.md.
% Separate items using '\ \ $|$ \ \ '. Wrap email in \href{mailto:...}{...}.
\begin{tcolorbox}[
  colback=bannerblue,
  colframe=bannerblue,
  arc=0mm,
  boxrule=0pt,
  top=5pt,
  bottom=5pt,
  left=10pt,
  right=10pt,
  halign=center
]
  \small \bfseries \color{textdark}
  [INSTRUCTION: City, Country] \ \ $|$ \ \ [INSTRUCTION: Phone Number] \ \ $|$ \ \ \href{mailto:[INSTRUCTION: Email Address]}{[INSTRUCTION: Email Address]}
\end{tcolorbox}

\vspace{5pt}

%==================================================
% NAME & TITLE
%==================================================
% INSTRUCTION FOR AGENT:
% 1. Line 1: Candidate's full legal name in UPPERCASE bold.
% 2. Line 2: Target job title / primary professional identity in uppercase darkblue,
%    optionally followed by ' $|$ ' and secondary specialization matching the target job posting.
\noindent
{\LARGE \textbf{[INSTRUCTION: CANDIDATE FULL NAME IN UPPERCASE]}}\\[3pt]
{\small \textbf{\color{darkblue}[INSTRUCTION: TARGET JOB TITLE $|$ SECONDARY SPECIALIZATION / DOMAIN]}}

\vspace{6pt}

%==================================================
% WEBSITES, PORTFOLIOS, PROFILES
%==================================================
% INSTRUCTION FOR AGENT:
% Extract 2-3 primary professional profile links (LinkedIn, GitHub, Portfolio) from Personal_Information.md.
% Use an itemize environment. Format each item as: \href{https://...}{display url without https://}.
\sectionline
\cvblock{WEBSITES,\\PORTFOLIOS,\\PROFILES}{%
  \begin{itemize}
    \item \href{[INSTRUCTION: Full URL to LinkedIn Profile]}{[INSTRUCTION: Clean Display URL, e.g., linkedin.com/in/username]}
    \item \href{[INSTRUCTION: Full URL to GitHub Profile]}{[INSTRUCTION: Clean Display URL, e.g., github.com/username]}
    % INSTRUCTION: Optional additional portfolio, website, or Google Scholar link if relevant:
    % \item \href{[INSTRUCTION: Full URL]}{[INSTRUCTION: Clean Display URL]}
  \end{itemize}
}

%==================================================
% PROFESSIONAL SUMMARY
%==================================================
% INSTRUCTION FOR AGENT:
% 1. Synthesize a powerful 3-5 sentence summary (~80-110 words) drawing from Professional_Summary.md,
%    CareerVision.md, and Keywords.md.
% 2. Tailor directly to the target job description, highlighting relevant years of experience,
%    key domain expertise, core technologies (bolded using \textbf{...}), problem-solving capabilities,
%    and language proficiencies.
% 3. Maintain absolute factual truth; do NOT invent qualifications or experience.
\sectionline
\cvblock{PROFESSIONAL\\SUMMARY}{%
  [INSTRUCTION: Write a tailored 3-5 sentence professional summary here. Highlight core specialization, years of experience, bolded key technologies matching target job, major domains of expertise, key achievements, and language proficiencies (e.g., English (Fluent / C1) and Spanish (Native / C2)). Ensure tone is confident, professional, and commercially grounded.]
}

%==================================================
% WORK EXPERIENCE
%==================================================
% INSTRUCTION FOR AGENT:
% 1. Source positions from Work_Experience/ and metrics from Quantifiable_Results.md.
% 2. Order positions in STRICT REVERSE CHRONOLOGICAL ORDER.
% 3. CRITICAL: NEVER alter, round, or fabricate dates, job titles, or company names.
% 4. Select the most relevant 3-5 roles. Emphasize roles closely aligned to the target job.
% 5. The FIRST role must be placed directly inside the \cvblock{WORK\\EXPERIENCE}{...} macro.
% 6. For each subsequent role, use the \jobentry{...} macro.
% 7. Bullet point guidelines:
%    - 3-5 bullets for primary/recent roles; 2-3 bullets for earlier roles.
%    - Structure: Action Verb + Context/Problem + Specific Solution/Tools + Quantifiable Impact.
%    - Bold critical technologies, tools, and methodologies (\textbf{Tool}).
%    - Escape LaTeX characters (\&, \%, \_).
\sectionline
\cvblock{WORK\\EXPERIENCE}{%
  \textbf{[INSTRUCTION: Most Recent Job Title]}\hfill \textbf{[INSTRUCTION: Start Month Year -- Present / End Month Year]}\\
  \emph{[INSTRUCTION: Company Name - City, Country (or Remote)]}
  \begin{itemize}
    \item {[INSTRUCTION: High-impact accomplishment bullet starting with strong action verb, detailing architecture, development, or methodology, bolding key tools with \textbf{Tech}, and including quantified result/metrics].}
    \item {[INSTRUCTION: Business value / operational impact bullet detailing how solutions solved user/stakeholder problems, saved time, cut costs, or scaled systems].}
    \item {[INSTRUCTION: End-to-end delivery bullet covering tech stack, deployment, infrastructure, or data/ML pipelines (\textbf{Technologies Used})].}
    \item {[INSTRUCTION: Leadership, cross-functional collaboration, mentoring, or engineering standards bullet].}
  \end{itemize}
}

% INSTRUCTION: Use \jobentry{...} for each additional relevant position (typically 2-4 more roles).
\jobentry{%
  \textbf{[INSTRUCTION: Second Most Recent Job Title]}\hfill \textbf{[INSTRUCTION: Start Month Year -- End Month Year]}\\
  \emph{[INSTRUCTION: Company Name - City, Country (or Remote)]}
  \begin{itemize}
    \item {[INSTRUCTION: High-impact tailored bullet with action verb, technologies (\textbf{Tech}), and measurable metric].}
    \item {[INSTRUCTION: Technical implementation / architecture / analysis bullet with quantified outcome].}
    \item {[INSTRUCTION: Collaboration or process automation bullet demonstrating operational value].}
  \end{itemize}
}

\jobentry{%
  \textbf{[INSTRUCTION: Third Most Recent Job Title]}\hfill \textbf{[INSTRUCTION: Start Month Year -- End Month Year]}\\
  \emph{[INSTRUCTION: Company Name - City, Country (or Remote)]}
  \begin{itemize}
    \item {[INSTRUCTION: High-impact tailored bullet highlighting technologies relevant to target role].}
    \item {[INSTRUCTION: Key accomplishment or system improvement with concrete outcome].}
  \end{itemize}
}

% INSTRUCTION: Add further \jobentry blocks if relevant and space permits, or omit earlier less-relevant roles to fit page limit.

%==================================================
% FEATURED PROJECTS
%==================================================
% INSTRUCTION FOR AGENT:
% 1. Source from Projects.md. Select 2-3 projects most relevant to the target job description.
% 2. Structure each item:
%    \item \textbf{[Project Name]:} [1-2 sentences explaining the project goal, system architecture,
%          and concrete outcome/value] (\textbf{[Tech 1]}, \textbf{[Tech 2]}, \textbf{[Tech 3]}).
% 3. Bold key technologies matching keywords from the job posting.
\sectionline
\cvblock{FEATURED\\PROJECTS}{%
  \begin{itemize}
    \item \textbf{[INSTRUCTION: Featured Project 1 Name]:} [INSTRUCTION: 1-2 sentence description explaining problem solved, technical approach, and quantifiable result/deliverable] (\textbf{[INSTRUCTION: Primary Technologies/Frameworks]}).
    \item \textbf{[INSTRUCTION: Featured Project 2 Name]:} [INSTRUCTION: 1-2 sentence description explaining problem solved, technical approach, and quantifiable result/deliverable] (\textbf{[INSTRUCTION: Primary Technologies/Frameworks]}).
    \item \textbf{[INSTRUCTION: Featured Project 3 Name]:} [INSTRUCTION: 1-2 sentence description explaining problem solved, technical approach, and quantifiable result/deliverable] (\textbf{[INSTRUCTION: Primary Technologies/Frameworks]}).
  \end{itemize}
}

%==================================================
% EDUCATION
%==================================================
% INSTRUCTION FOR AGENT:
% 1. Source degrees from Education.md in REVERSE CHRONOLOGICAL ORDER.
% 2. CRITICAL: NEVER alter degree titles, dates, institutions, or locations.
% 3. Format each degree entry with:
%    \textbf{[Degree Title]}\hfill \textbf{[Start Date -- End Date / Present]}\\
%    \emph{[Institution Name] - [City, Country]}\\
%    {\small \emph{Key Coursework:} [Course 1, Course 2, Course 3, ... (tailored to job)].}
% 4. Include honors, GPA, or academic exchange scholarships (e.g., ERASMUS+) where applicable.
% 5. Separate entries with \vspace{3pt}.
\sectionline
\cvblock{EDUCATION}{%
  \textbf{[INSTRUCTION: Master's / Highest Degree Title]}\hfill \textbf{[INSTRUCTION: Start Month Year -- Graduation Month Year / Present]}\\
  \emph{[INSTRUCTION: University Name - City, Country]}\\
  {\small \emph{Key Coursework:} [INSTRUCTION: 4-6 coursework topics most relevant to target job description.]}

  \vspace{3pt}

  \textbf{[INSTRUCTION: Bachelor's / Undergraduate Degree Title]}\hfill \textbf{[INSTRUCTION: Start Month Year -- Graduation Month Year]}\\
  \emph{[INSTRUCTION: University Name - City, Country]}\\
  {\small \emph{Key Coursework:} [INSTRUCTION: 4-6 relevant coursework topics.] [INSTRUCTION: GPA if noteworthy / required.]}

  % INSTRUCTION: Optional academic exchange, honors, or scholarships:
  \vspace{3pt}

  \textbf{[INSTRUCTION: Academic Grant / Exchange Program / Honor (if applicable)]}\hfill \textbf{[INSTRUCTION: Year / Duration]}\\
  {\small [INSTRUCTION: Brief 1-line description of research focus, scholarship prestige, or academic accomplishments.]}
}

%==================================================
% CERTIFICATIONS AND RECOGNITION
%==================================================
% INSTRUCTION FOR AGENT:
% 1. Source from Certifications.md, Publications_Talks.md, and Achievements.md.
% 2. Prioritize industry-recognized certifications and professional honors relevant to target job.
% 3. Format as a dense, comma-separated paragraph to conserve vertical space.
% 4. Include peer-review activities, patents, or publications if relevant to role.
\sectionline
\cvblock{CERTIFICATIONS\\AND HONORS}{%
  [INSTRUCTION: Comma-separated list of relevant certifications, credentialing bodies (e.g., AWS, edX, Harvard, deeplearning.ai), academic grants, and peer-review / editorial honors tailored to the target position.]
}

%==================================================
% TECHNICAL SKILLS
%==================================================
% INSTRUCTION FOR AGENT:
% 1. Source from Skills.md and Technologies.md.
% 2. Align selection with the job posting requirements and ATS keywords.
% 3. Group into 2-4 coherent skill categories (e.g., Applied AI \& ML, Software \& Cloud, Data \& Analytics).
% 4. Format: \textbf{[Category Name]:} Skill1, Skill2, Skill3 \ \ $|$ \ \ \textbf{[Next Category]:} ...
% 5. Only include skills and tools grounded in the candidate's source documents.
\sectionline
\cvblock{TECHNICAL SKILLS}{%
  \textbf{[INSTRUCTION: Category 1 (e.g., Primary Specialty / Applied AI \& ML)]:} [INSTRUCTION: Comma-separated list of relevant tools, libraries, models, and frameworks] \ \ $|$ \ \ \textbf{[INSTRUCTION: Category 2 (e.g., Software Engineering \& Cloud)]:} [INSTRUCTION: Comma-separated list of languages, cloud services, and DevOps tools] \ \ $|$ \ \ \textbf{[INSTRUCTION: Category 3 (e.g., Data \& Analytics)]:} [INSTRUCTION: Comma-separated list of databases, ETL, BI, and data processing tools]
}

%==================================================
% CORE COMPETENCIES
%==================================================
% INSTRUCTION FOR AGENT:
% 1. Source from Skills.md, Keywords.md, and Leadership.md.
% 2. Select 8-10 high-value functional competencies, architectural strengths, and leadership attributes matching the target job.
% 3. Present in a two-column layout using two minipages (0.48\linewidth each) with itemize lists.
% 4. Include language proficiencies in the final bullet (e.g., English (Fluent C1) \& Spanish (Native C2)).
\sectionline
\cvblock{CORE\\COMPETENCIES}{%
  \begin{minipage}[t]{0.48\linewidth}
    \begin{itemize}
      \item {[INSTRUCTION: Core competency 1 tailored to target role]}
      \item {[INSTRUCTION: Core competency 2 tailored to target role]}
      \item {[INSTRUCTION: Core competency 3 tailored to target role]}
      \item {[INSTRUCTION: Core competency 4 tailored to target role]}
      \item {[INSTRUCTION: Core competency 5 tailored to target role]}
    \end{itemize}
  \end{minipage}%
  \hfill%
  \begin{minipage}[t]{0.48\linewidth}
    \begin{itemize}
      \item {[INSTRUCTION: Core competency 6 tailored to target role]}
      \item {[INSTRUCTION: Core competency 7 tailored to target role]}
      \item {[INSTRUCTION: Core competency 8 tailored to target role]}
      \item {[INSTRUCTION: Core competency 9 tailored to target role]}
      \item {[INSTRUCTION: Language proficiencies from Personal\_Information.md, e.g., English (Fluent C1) \& Spanish (Native C2)]}
    \end{itemize}
  \end{minipage}
}

\end{document}
